% This file is based on the "sig-alternate.tex" V1.9 April 2009
% This file should be compiled with V2.4 of "sig-alternate.cls" April 2009

\documentclass{sig-alternate}

\usepackage{url}
\usepackage{color}
\usepackage{enumerate}
\usepackage{balance}
\permission{}
\CopyrightYear{2013}
%\crdata{0-00000-00-0/00/00}
\begin{document}

\title{Holographic Pong}
\numberofauthors{1}
\author{ 
\alignauthor
Piter Garcia,
Vasudev Prasad,
Jorge Leon
}
\date{25 August 2013}
\maketitle
\begin{abstract}
Holography is a technique used to display objects or scenes in three dimensions. These 3-dimensional images, or holograms, can be seen with the unassisted eye and they are very similar to how humans see the actual environment that surround them. Unfortunately, this technique has not been very successful, principally due to the computational complexity of calculating diffraction patterns in real time and poor quality of the resultant images. In our project, our goal is to overcome these problems to enable a high-quality real-time holographic projection of a pong game that we are developing. In order to approach our goal, we are making use of a holographic stereo-graphic technique and a polymer material as the medium to demonstrate a holographic display of our pong game. A projector is used to project the holographic image or multicolored holographic 3D images produced by our code implementation. The 3D feeling is given when taking the image projected from one location to another location where the hologram is projected in real-time as a 3D display. Our project could easily bring future applications in the entertainment area, and other areas as well, with only some further improvements.
\end{abstract}

\section{Overview}
\label{overview}


\section{Design}
\label{design considerations}


\section{Architecture}
\label{architecture}


\section{Implementation}
\label{implementation}

\section{Lessons Learned}
\label{mistakes}


\section{Current Status and Future Work}
\label{current status}


\subsection{Tables, Figures}
\section{Citations/References}


\bibliographystyle{abbrv}
\bibliography{report}
% You must have a proper ".bib" file
%  and remember to run:
% latex bibtex latex latex
% to resolve all references
\balance
\end{document}








